% Generated draft: requires mathematical and source review.
\documentclass[11pt]{article}
\usepackage[margin=25mm]{geometry}
\usepackage{fontspec}
\setmainfont{Latin Modern Roman}
\usepackage{amsmath,amssymb,mathtools,unicode-math}
\setmathfont{Latin Modern Math}
\usepackage{enumitem,needspace,xurl,hyperref}
\setlength{\parindent}{0pt}
\setlength{\parskip}{5pt}
\setlength{\emergencystretch}{3em}
\newcommand{\N}{\mathbb{N}}
\newcommand{\Z}{\mathbb{Z}}
\newcommand{\Q}{\mathbb{Q}}
\newcommand{\R}{\mathbb{R}}
\newcommand{\C}{\mathbb{C}}
\newcommand{\F}{\mathbb{F}}
\newcommand{\A}{\mathbb{A}}
\newcommand{\PP}{\mathbb{P}}
\newcommand{\E}{\mathbb{E}}
\newcommand{\eps}{\varepsilon}
\newcommand{\dd}{\,\mathrm{d}}
\begin{document}
\section*{UnsolvedMath ID: 701182}
\textbf{Source problem code:} catalogue entry
\subsection{Definitions and mathematical statement}
For a group \(H\) and a subset \(S\subseteq H\), define the normal closure of \(S\) in \(H\) by
\[
\langle\!\langle S\rangle\!\rangle_H
=\bigl\langle hsh^{-1}:h\in H,\ s\in S\bigr\rangle.
\]
The weight of \(H\) is the cardinal
\[
w(H)=\min\bigl\{|S|:S\subseteq H,\ \langle\!\langle S\rangle\!\rangle_H=H\bigr\}.
\]
This definition permits infinite weight.

Let \(n\geq 1\) be a finite integer, let \(G_1,\ldots,G_n\) be nontrivial groups, and let
\[
G=G_1*\cdots*G_n
\]
be their free product. J. Howie's problem asks the following questions.
\begin{enumerate}
\item[(a)] For every such choice of \(n\) and factors, is
\[
w(G)\geq \left\lceil\frac n2\right\rceil?
\]
\item[(b)] Does the same inequality hold for every such choice when each \(G_i\) is cyclic? The cyclic factors may be finite or infinite, but must be nontrivial.
\item[(c)] For every three nontrivial groups \(G_1,G_2,G_3\), is
\[
w(G_1*G_2*G_3)\geq 2?
\]
\end{enumerate}
The bound in (a) and (b) expresses the source's inequality \(w(G)\geq n/2\): no subset with fewer than \(\lceil n/2\rceil\) elements should normally generate \(G\).

The supplied record states that Archive Problem 5.53, the special case \(n=3\) with all three factors cyclic, has been answered affirmatively. This is a supplied historical qualification, not an independently verified result here; no separate solution reference is supplied. The registry labels Problem 15.94 as open, which is likewise a supplied status rather than a new literature assessment.
\subsection{Short English statement}
The weight of a group is the least number of elements needed to generate it normally. For a free product of \(n\) nontrivial groups, must the weight be at least \(\lceil n/2\rceil\)? The problem also asks this with cyclic factors and, separately, with three arbitrary nontrivial factors. The source reports an affirmative answer for three cyclic factors.
\subsection{Sources}
\begin{itemize}
\item Local UnsolvedMath registry, permanent ID \(701182\), exact source code \texttt{catalogue entry}, title \emph{Kourovka Notebook Problem 15.94}; supplied complete source record.
\item J. Howie, proposer, Problem 15.94 in \emph{The Kourovka Notebook (Unsolved Problems in Group Theory)}, 15th issue (2002). The supplied record also identifies the 21st issue edition (2026), arXiv:1401.0300v48. The collection metadata names E. I. Khukhro and V. D. Mazurov as editors.
\item The supplied statement refers to Archive Problem 5.53 and reports that its three-cyclic-factor case has been answered affirmatively. No separate solution citation is provided. These source and status descriptions have not been independently verified here.
\end{itemize}
\end{document}
